\section{结构因果模型把干预写成机制替换}
\label{sec:13-Machine-Learning/12-Causal-Inference/02}

一个电商团队想知道推送优惠券能带来多少增量成交。历史数据里，收到券的用户成交率是 \(31\%\)，没收到券的是 \(12\%\)。差 \(19\) 个百分点，看着很可观。但发券规则是运营写的：近三十天浏览次数多的用户优先发。而浏览次数多的用户，本来成交率就高。

团队的第一反应是``控制掉浏览次数''，于是在浏览分层内比较——这一步方向对了，可要说清为什么这样算对、什么时候不对，得先有一个能写下``发券规则被替换了''这件事的语言。条件概率写不出它：条件化只是筛人，而发券规则改成``全员发''是把一条机制换掉了。

\subsection{一条机制就是一个赋值方程}

结构因果模型给每个内生变量写一条生成方程，方程右边列出它的直接原因和一团未被进一步解释的外生扰动。上面的场景写成

\[
\begin{aligned}
C &= f_C(U_C),\\
X &= f_X(C,U_X),\\
Y &= f_Y(X,C,U_Y),
\end{aligned}
\]

其中 \(C\) 是处理前协变量（浏览次数），\(X\) 是是否发券，\(Y\) 是是否成交。这里的等号是赋值，不是等式两边可以随便移项的那种相等：\(f_X\) 说的是``\(X\) 由 \(C\) 和 \(U_X\) 产生''，把它反过来写成 \(C\) 关于 \(X\) 的函数就改变了模型的含义。箭头由函数依赖读出，得到 \(C\to X\)、\(C\to Y\)、\(X\to Y\)。

比起一个联合分布，这套方程多说了一句话：机制是模块化的。每条方程代表一个可以被单独换掉的物理或制度过程，换掉其中一条，其余各条照旧运转。这句话不是从数据里学来的，是建模者的断言，而且它正是干预演算全部的依据。

执行 \(\operatorname{do}(X=x)\) 就是把 \(X=f_X(C,U_X)\) 整条删掉，换成常数方程

\[
X:=x,
\]

其余方程一字不动。\(C\) 仍按 \(f_C\) 产生，\(Y\) 仍按 \(f_Y\) 产生，只是 \(f_Y\) 的第一个参数现在被外部按住。运营改成全员发券，就是把``浏览多的优先''这条规则替换掉，而用户看到券之后要不要下单的那套心理过程假定没变。这个假定一旦不成立——比如全员发券让券显得廉价、转化机制随之改变——替换就是错的，但至少它被写在了纸面上，而不是藏在回归系数背后。

\subsection{删一个因子和重新归一化是两种算法}

把模块化落到分布上，就得到干预分布的计算规则。观察联合分布按图分解为

\[
p(v_1,\ldots,v_d)=\prod_{j=1}^{d}p\bigl(v_j\given v_{\operatorname{pa}(j)}\bigr),
\]

干预 \(X\) 之后，\(X\) 自己那个因子被删除，其余因子原样保留并把 \(X\) 固定在 \(x\)：

\[
p\bigl(v_{-X}\given\operatorname{do}(X=x)\bigr)
=\prod_{j:\,V_j\ne X}p\bigl(v_j\given v_{\operatorname{pa}(j)}\bigr)\bigg|_{X=x}.
\]

这叫截断分解。它不是任何一个 Bayesian 网络自动具备的性质——同一个联合分布可以有多种因子分解，删哪一个才对，取决于哪张图被赋予了机制含义。

对照一下条件化。在三变量的例子里 \(p(c,x,y)=p(c)p(x\given c)p(y\given x,c)\)，两种操作的差别可以摆在一起看：

\[
\begin{aligned}
p(y\given X=x)&=\sum_c p(y\given x,c)\,p(c\given x),\\[2pt]
p\bigl(y\given\operatorname{do}(X=x)\bigr)&=\sum_c p(y\given x,c)\,p(c).
\end{aligned}
\]

同样是对 \(c\) 求和，同样用 \(p(y\given x,c)\) 做被加权项，差别全在权重。条件化把 \(p(x\given c)\) 留着，靠 Bayes 反过来变成 \(p(c\given x)\)——发了券的那群人里浏览多的占比本来就高，这个权重把选择规则带进了答案。干预把 \(p(x\given c)\) 删掉，剩下的权重是目标总体自己的 \(p(c)\)。第二个式子就是标准化公式，右边每一项都能从观察数据估出来，所以在这张图下干预分布可识别。

\begin{center}
\begin{tikzpicture}[>=stealth,line width=0.7pt]

% ---- left: observational graph ----
\begin{scope}[shift={(0,0)}]
  \draw (0,1.3) circle (0.34); \node at (0,1.3) {\(C\)};
  \draw (-1.3,-0.3) circle (0.34); \node at (-1.3,-0.3) {\(X\)};
  \draw (1.3,-0.3) circle (0.34); \node at (1.3,-0.3) {\(Y\)};
  \draw[->] (-0.26,1.07) -- (-1.04,-0.07);
  \draw[->] (0.26,1.07) -- (1.04,-0.07);
  \draw[->] (-0.96,-0.3) -- (0.96,-0.3);
  \node at (0,-1.25) {\footnotesize \(p(c)\,p(x\given c)\,p(y\given x,c)\)};
  \node at (0,-1.95) {\footnotesize 观察：权重是 \(p(c\given x)\)};
\end{scope}

% ---- arrow ----
\draw[->,line width=1pt] (2.5,-0.3) -- (3.9,-0.3);
\node at (3.2,0.1) {\footnotesize \(\operatorname{do}(X=x)\)};

% ---- right: mutilated graph ----
\begin{scope}[shift={(6.4,0)}]
  \draw (0,1.3) circle (0.34); \node at (0,1.3) {\(C\)};
  \draw[line width=1.1pt] (-1.3,-0.3) circle (0.34); \node at (-1.3,-0.3) {\(x\)};
  \draw (1.3,-0.3) circle (0.34); \node at (1.3,-0.3) {\(Y\)};
  \draw[->,dashed,gray!60] (-0.26,1.07) -- (-1.04,-0.07);
  \draw[line width=1.1pt] (-0.85,0.72) -- (-0.35,0.28);
  \draw[line width=1.1pt] (-0.85,0.28) -- (-0.35,0.72);
  \draw[->] (0.26,1.07) -- (1.04,-0.07);
  \draw[->] (-0.96,-0.3) -- (0.96,-0.3);
  \node at (0,-1.25) {\footnotesize \(p(c)\,\phantom{p(x\given c)}\,p(y\given x,c)\)};
  \draw[line width=0.9pt] (-0.60,-1.25) -- (0.14,-1.25);
  \node at (0,-1.95) {\footnotesize 干预：权重是 \(p(c)\)};
\end{scope}

\end{tikzpicture}
\end{center}

\emph{干预剪掉指向 \(X\) 的箭头，等价于从乘积里划掉 \(p(x\given c)\) 这一个因子；条件化一个因子也不删，只是把整个乘积重新归一化。}

图右边那道横线是这一节的全部内容：被划掉的因子恰好是编码了``谁更容易拿到券''的那一项。条件化留着它，于是答案里混着运营的发券规则；干预删掉它，答案只反映券本身的作用。回到开头的数字，团队要算的是右式，而 \(19\) 个百分点是左式。

连续的 \(C\) 把求和换成积分。\(C\) 维数一高，直接分层往往每层没几个样本，后面会改用回归或倾向得分来估同一个目标——换的是计算方式，识别公式的前提一条没少。

\subsection{模块化是假设，它有失效的具体方式}

写下截断分解的那一刻，你已经承诺了几件事，值得逐条对账。

其一，被干预变量之外的方程不随干预改变。全员发券若使券的稀缺感消失、\(f_Y\) 随之变形，替换就失效了。其二，干预本身有明确定义。``让某人减肥十公斤''不是一条机制——节食、运动、患病各自对应不同的 \(f\)，它们的下游后果不同，\(\operatorname{do}\) 记号掩盖了这个歧义。其三，单位之间互不影响。给一部分用户发券若挤占了另一部分用户的库存，一个人的 \(X\) 会进入另一个人的 \(Y\) 方程，模型写的那套单机制方程就不够用了。其四，图里没有漏掉共同原因。漏一个，截断分解的每一步照样能写下来，只是写下来的东西不再对应任何真实干预。

这些条件里没有一条能靠增大样本量满足。它们是结构性的，只能靠设计、领域知识或额外的实验来支撑。因果图的价值也在这里——它把这份清单摆在明处，让识别推导逐条可审，箭头画着直观只是顺带的好处。

至于图本身从哪来：观察数据能定到的那一层非常有限，两张蕴涵相同条件独立关系的图任何检验都分不开，这条界限的代数刻画见\hyperref[sec:12-Real-Analysis-Support/10-Essential-Structures/04]{``因果结构的代数形式''}一节。本章接受图是给定的，转而追问一个更实际的问题：图在手上时，究竟该控制哪些变量。下一节会说明，答案既不是``全都控制''，也不是``只控制看起来相关的''。
